This internship project set out to design and build a real-time anomaly
detection platform for branch teller operations at Amen Bank, covering
four operation types: retrait, versement, virement, and remise de
chèques. The result is a working end-to-end system that ingests
transaction events, computes behavioral features, scores them with two
complementary deep learning models, applies regulatory rules, and
surfaces explainable alerts to branch supervisors through an interactive
dashboard.

The platform monitors two independent actors --- the client initiating
the transaction and the employee processing it --- reflecting the
dual-actor nature of branch teller operations where anomalies can
originate from either side. A streaming architecture built on RedPanda
chains four microservices through message topics, supported by a
hot/cold data layer (Redis and PostgreSQL) and a batch pipeline for
nightly profile reconstruction. The Autoencoder detects point anomalies
in individual transactions while the LSTM identifies sequential
deviations in client behavior over time, with a sigmoid-weighted fusion
mechanism that adapts to account maturity. Every alert carries a
human-readable SHAP explanation generated by deterministic templates,
enabling supervisors to understand and act on the system's output
without requiring machine learning expertise.

Beyond the detection pipeline, the project invested heavily in the
infrastructure that makes a machine learning system maintainable. MLflow
tracks every training experiment. An automated promotion gate ensures
that only models matching or exceeding current production performance
are deployed. Three independent retraining triggers --- a weekly
schedule, a Prometheus-driven drift detector, and a supervisor feedback
divergence monitor --- close the loop between model serving and model
improvement. Structured JSON logging, Prometheus metrics, and Grafana
dashboards provide the observability layer needed to diagnose issues
without inspecting individual events.

The system was validated on a synthetic dataset of 243,000 events
generated from five client archetypes calibrated against Amen Bank's
published data. The validation confirms that the pipeline works
end-to-end --- from event ingestion to supervisor alert --- and that
the retraining subsystem can detect drift and trigger model updates
autonomously. The synthetic data caveat is explicit: these results
demonstrate engineering readiness, not production-grade detection
accuracy. The architecture is designed so that replacing synthetic
data with real transactions requires no code changes --- the same
feature computation, the same models, and the same retraining
pipeline apply without modification.

This project provided the opportunity to work across the full spectrum
of a machine learning system: from data engineering and model training
to streaming infrastructure, MLOps, and front-end delivery. The most
valuable lesson was that building a model is a small fraction of the
effort --- the majority of the engineering work lies in the
infrastructure that ensures the model remains correct, observable, and
improvable over time.